\documentclass[11pt,letterpaper]{article}
\usepackage{fontspec}
\usepackage{tocloft}
\usepackage{multicol}
\usepackage{fancyhdr}
\usepackage{nicefrac}
\usepackage[left=1.4in,right=1.5in,top=1in,bottom=1in]{geometry}
\setmainfont[Scale=1.4,AutoFakeBold=1.5,AutoFakeSlant=0.3]{Doves Type}

% Prevent page breaks within list items
\widowpenalty=10000
\clubpenalty=10000
\interlinepenalty=500

\title{Gluten-Free Tuna Mac •}
\author{}
\date{}

% Configure fancy header
\pagestyle{fancy}
\fancyhf{}
\fancyhead[C]{\textit{Gluten-Free Tuna Mac}}
\fancyfoot[C]{\thepage}
\renewcommand{\headrulewidth}{0pt}
\renewcommand{\footrulewidth}{0pt}

\begin{document}

\maketitle
\thispagestyle{empty}

\begin{quote}\textit{A quick stovetop comfort dish that upgrades classic boxed mac and cheese into a savory weeknight supper. Warm herbs and a touch of acid cut through the rich, velvety cheese sauce so the pasta and tender flakes of tuna taste bright rather than heavy. Serves 3--4.}\end{quote}

\section*{Ingredients}
\setlength{\columnsep}{20pt}
\begin{multicols}{2}
\noindent
    Whole milk \dotfill 2~cups \\
    Dried onion flakes \dotfill 1~tsp. \\
    Dried tarragon \dotfill \nicefrac{1}{2}~tsp. \\
    Ground nutmeg \dotfill \nicefrac{1}{8}~tsp. \\
    Salted butter, divided \dotfill 3~Tbsp. \\
    White rice flour \dotfill 2~Tbsp. \\
    Kosher salt, divided \dotfill 4~tsp. \\
    White or black pepper \dotfill \nicefrac{1}{2}~tsp. \\
    Ground mustard \dotfill \nicefrac{1}{2}~tsp. \\

    \columnbreak
    \noindent
    Ground fennel \dotfill \nicefrac{1}{8}~tsp. \\
    American cheese \dotfill 1~slice (\nicefrac{3}{4}~oz.) \\
    Cheddar or Gruyère cheese \dotfill 4--6~oz. \\
    Gluten-free elbow pasta \dotfill 6--8~oz. \\
    Yellow mustard \dotfill 1~tsp. \\
    White wine vinegar \dotfill 1~tsp. \\
    Canned tuna \dotfill 1~can (6~oz.) \\
    Dried parsley \dotfill 2~tsp. \\
\end{multicols}

\section*{Sauce}

\noindent
Bring a large pot of \textbf{water} to a boil ---
Shred \textbf{sharp~cheddar} or \textbf{Gruyère~cheese}

\begin{enumerate}
    \item Warm 2~cups~\textbf{whole~milk}, 1~tsp.~\textbf{dried~onion~flakes}, \nicefrac{1}{2}~tsp.~\textbf{dried~tarragon}, and \nicefrac{1}{8}~tsp.~\textbf{ground~nutmeg} in the microwave on \textit{high} for \textit{2--3~minutes} until steaming hot (\textit{165°F}) and the \textbf{onion~flakes} and \textbf{tarragon} soften.

    \item Melt 2~Tbsp.~\textbf{salted~butter} in a medium saucepan over \textit{medium~heat} for \textit{30--45~seconds} until foamy. Whisk in 2~Tbsp.~\textbf{white~rice~flour} and cook for \textit{1~minute} until smooth and bubbling.

    \item Stir in 1~tsp.~\textbf{kosher~salt}, \nicefrac{1}{2}~tsp.~\textbf{ground~pepper}, \nicefrac{1}{2}~tsp.~\textbf{ground~mustard}, and \nicefrac{1}{8}~tsp.~\textbf{ground~fennel}. Cook, whisking, for \textit{1~minute} until foamy, pale gold, and fragrant.

    \item Whisk in the warm \textbf{milk} a splash at a time so the roux catches and stays smooth before adding more liquid.

    \item Bring to a low simmer over \textit{medium-low~heat} and cook, whisking often, for \textit{3--5~minutes} until smooth and thick enough to coat the back of a spoon (\textit{185--195°F}). Continue simmering for \textit{1--2~minutes} if any floury grit remains.

    \item Reduce to \textit{low~heat}. Whisk in 1~slice~\textbf{American~cheese} until melted, about \textit{30--60~seconds}. Add the shredded \textbf{cheddar} or \textbf{Gruyère~cheese} a handful at a time, stirring until each handful melts smoothly before adding the next, about \textit{2--3~minutes}.

    \item Stir in the remaining 1~Tbsp.~\textbf{salted~butter} until melted. Taste and adjust \textbf{salt}. Keep on the \textit{lowest~heat} (below \textit{180°F}) while the \textbf{pasta} cooks, stirring occasionally.
\end{enumerate}

\section*{Finish}

\noindent
Drain \textbf{tuna}

\begin{enumerate}
    \item Add 1~Tbsp.~\textbf{kosher~salt} and 6--8~oz.~\textbf{gluten-free~elbow~pasta} to the boiling \textbf{water}. Boil over \textit{medium-high~heat} for \textit{6--8~minutes}, stirring occasionally, until al dente (tender with a slight bite). Drain well and return the \textbf{elbow~pasta} to the pot.

    \item Remove the sauce from the heat and stir in 1~tsp.~\textbf{yellow~mustard} and 1~tsp.~\textbf{white~wine~vinegar}.

    \item Pour the cheese sauce over the \textbf{elbow~pasta} in the pot, then add the drained \textbf{tuna} and 2~tsp.~\textbf{dried~parsley}. Fold gently for \textit{1~minute} until coated and heated through. Taste for \textbf{salt} and serve hot.
\end{enumerate}

\newpage

{\footnotesize
\setlength{\columnsep}{20pt}
\setlength{\multicolsep}{6pt}
\begin{multicols}{2}
\setlength{\parindent}{0pt}
\setlength{\parskip}{2pt}
\setlength{\itemsep}{0pt}
\setlength{\parsep}{0pt}

\section*{Equipment Required}
\begin{itemize}
    \item Large pot (4--6~quart, for \textbf{pasta})
    \item Medium saucepan (3-quart, for sauce)
    \item 4-cup microwave-safe measuring cup
    \item Box grater
    \item Colander and can opener
    \item Whisk and silicone spatula
    \item Measuring cups and spoons
\end{itemize}

\section*{Hints and Notes}
\subsection*{Yield}
\begin{itemize}
    \item Serves 3--4 as a main dish
\end{itemize}

\subsection*{Mise en Place}
\begin{itemize}
    \item Start the \textbf{pasta} \textbf{water} heating first so it boils while the sauce simmers
    \item Shred \textbf{cheddar} or \textbf{Gruyère cheese} and drain \textbf{tuna} before starting the roux
    \item Warm the \textbf{whole milk}, \textbf{dried onion flakes}, \textbf{dried tarragon}, and \textbf{ground nutmeg} right before melting the \textbf{salted butter}
\end{itemize}

\subsection*{Ingredient Tips}
\begin{itemize}
    \item \textbf{White rice flour}: Use finely milled \textbf{white rice flour} so the sauce cooks up smooth
    \item \textbf{American cheese}: A single slice keeps the \textbf{cheddar} or \textbf{Gruyère} smooth and prevents the sauce from splitting
    \item \textbf{Cheddar} or \textbf{Gruyère}: Shred from a block (or blend both); pre-shredded \textbf{cheese} contains starches that make the sauce grainy
    \item \textbf{Gluten-free elbow pasta}: Corn-and-rice blends (such as Barilla or Jovial) hold their shape best. Use 6~oz. for a saucier dish or 8~oz. for a thicker coat
    \item \textbf{Canned tuna}: Drain well so excess \textbf{water} or \textbf{oil} does not thin the sauce
\end{itemize}

\subsection*{Preparation Tips}
\begin{itemize}
    \item Steeping the \textbf{dried onion flakes}, \textbf{dried tarragon}, and \textbf{ground nutmeg} in the warm \textbf{whole milk} softens the herbs and infuses the liquid without scorching
    \item Cooking the \textbf{roux spices} in the \textbf{salted butter} and \textbf{white rice flour} for \textit{1~minute} brings out their aroma without burning
    \item Add the warm \textbf{milk} gradually at the start so the roux catches into a smooth paste before thinning
    \item \textbf{White rice flour} benefits from a gentle simmer (\textit{185--195°F}) to lose its chalky texture and turn glossy
    \item Keep on \textit{low heat} once the \textbf{cheese} is added; boiling above \textit{190°F} can make the sauce grainy
    \item Stir in the \textbf{white wine vinegar} and \textbf{yellow mustard} right before combining with the \textbf{pasta} for a bright finish
    \item Fold the \textbf{tuna} gently at the end so the flakes stay intact
\end{itemize}

\subsection*{Make Ahead \& Storage}
\begin{itemize}
    \item Best served right away while the sauce is hot and creamy
    \item Leftovers keep \textit{3~days} refrigerated in an airtight container
    \item Reheat gently over \textit{low heat} or in the microwave at \textit{50\%~power} with a splash of \textbf{milk} to loosen the sauce
    \item Do not freeze; \textbf{gluten-free pasta} softens when thawed
\end{itemize}

\subsection*{Serving Suggestions}
\begin{itemize}
    \item Serve hot in bowls
    \item Top with crushed \textbf{potato chips} or toasted \textbf{gluten-free breadcrumbs} for crunch
    \item Pair with steamed \textbf{peas}, roasted \textbf{broccoli}, or a simple green salad
    \item Offer \textbf{hot sauce} or extra \textbf{black pepper} at the table
\end{itemize}

\end{multicols}
}

\end{document}
